\subsection{\acrlong{MDF}}
\label{subsec:modes_deffaillance}


Les \acrfull{MDF} constituent les causes techniques susceptibles de conduire aux \acrshort{ERPs} identifiés en
\acrlong{APR}. Contrairement aux \acrshort{ERP}, qui décrivent des situations finales dangereuses au niveau système, les
\acrlong{MDF} décrivent les anomalies, ruptures ou dysfonctionnements pouvant apparaître au sein des sous-systèmes du lanceur.\\

L’objectif de cette section est d’identifier, pour chaque ensemble fonctionnel de SP-02, les défaillances susceptibles
de conduire à un ou plusieurs \acrshort{ERPs}. Cette étape permet de caractériser l’origine potentielle des risques avant leur
évaluation détaillée dans l’\acrshort{AMDEC} (section~\ref{subsec:AMDEC}).

\subsubsection{Présentation des sous-systèmes}
\label{subsubsec:sous_sys}

Pour l’analyse des \acrshort{MDF}, les sous-systèmes du lanceur SP-02 sont regroupés en six familles principales :\\
\begin{itemize}
    \item \textbf{Éléments mécaniques statiques} : structures porteuses, coiffe, ailerons, bagues d’interface ou de liaison,
    interface entre étages.
    \item \textbf{Éléments mécaniques dynamiques} : aérofreins, trappes parachutes, mécanismes de séparation.
    \item \textbf{Éléments électroniques} : capteurs, cartes séquenceurs, alimentation, chaîne pyrotechnique.
    \item \textbf{Éléments logiciels} : logique d’état, traitements capteurs, conditions d’autorisation.
    \item \textbf{Facteurs humains (procéduraux)} : opérations d’assemblage, d’armement et de préparation du vol.
    \item \textbf{Facteurs externes (environnementaux)} : conditions météorologiques, environnement électromagnétique.\\
\end{itemize}

Les paragraphes suivants présentent les \acrshort{MDF} génériques associés à ces familles. Chaque mode est noté
\textbf{\acrshort{MDF}$i$} et peut concerner plusieurs sous-systèmes ou familles lorsque cela est pertinent.

\subsubsubsection*{Définition des \acrshort{MDF} génériques}

Les \acrlong{MDF} suivants ont été retenus :

\begin{itemize}
    \item \textbf{\acrshort{MDF}1 – Perte d’intégrité structurelle} \\
    Rupture, fissuration, délamination ou flambage d’un élément structurel (peau porteuse, aileron, coiffe, bague),
    entraînant une modification significative de la rigidité ou de la géométrie.\\

    \item \textbf{\acrshort{MDF}2 – Géométrie ou alignement non conforme} \\
    Défaut d’alignement entre étages, flèche excessive (\og banane\fg{} du tube), jeu mécanique anormal ou montage hors
    tolérance, pouvant perturber la stabilité ou la séparation.\\

    \item \textbf{\acrshort{MDF}3 – Blocage ou mouvement incomplet d’un mécanisme} \\
    Blocage total ou partiel d’un mécanisme mobile (aérofrein, trappe parachute, loquet, ressort), dû à des frottements, à une
    déformation ou à une pollution.\\

    \item \textbf{\acrshort{MDF}4 – Mouvement asymétrique ou non maîtrisé} \\
    Déploiement dissymétrique (un seul aérofrein, ouverture partielle d’une trappe, séparation non symétrique), entraînant un
    moment non désiré ou une transition de phase dégradée.\\

    \item \textbf{\acrshort{MDF}5 – Perte ou corruption de mesure} \\
    Capteur hors service, saturé ou perturbé (IMU, altimètre, GPS), câble capteur sectionné ou pincé, ou bruit électromagnétique
    induisant des mesures incohérentes.\\

    \item \textbf{\acrshort{MDF}6 – Défaut d’alimentation électrique} \\
    Perte totale ou partielle d’alimentation (batterie déchargée, chute de tension, connecteur mal engagé), rendant inopérant un
    ou plusieurs sous-systèmes électroniques.\\

    \item \textbf{\acrshort{MDF}7 – Défaut de la chaîne pyrotechnique} \\
    Court-circuit, isolation insuffisante, impulsion parasite ou connectique inadaptée sur la ligne pyrotechnique du second étage
    ou des systèmes de récupération, pouvant conduire à un allumage intempestif ou à un non-allumage.\\

    \item \textbf{\acrshort{MDF}8 – Erreur de logique d’état ou de séquencement} \\
    Défaut de machine d’état (transition illégale, état bloqué), erreur de timer ou d’ordonnancement des événements (séparation,
    allumage, ouverture parachute).\\

    \item \textbf{\acrshort{MDF}9 – Interprétation erronée des conditions d’autorisation} \\
    Mauvaise prise en compte des conditions de séparation, d’attitude ou de fenêtre temporelle, conduisant à autoriser ou
    interdire à tort une action critique.\\

    \item \textbf{\acrshort{MDF}10 – Erreur procédurale ou d’assemblage} \\
    Erreur humaine lors du montage, du câblage, de l’armement ou du rangement (parachute mal conditionné, connecteur inversé,
    coiffe mal fixée, shunt mal géré).\\

    \item \textbf{\acrshort{MDF}11 – Préparation ou configuration non conforme} \\
    Configuration logicielle ou matérielle incorrecte (paramètres non mis à jour, seuils mal réglés, configuration d’essai
    laissée en place), ou non-respect d’une étape de préparation.\\

    \item \textbf{\acrshort{MDF}12 – Conditions environnementales hors domaine} \\
    Conditions extérieures (vent fort, pluie, humidité, température, perturbations électromagnétiques) dépassant le domaine pris
    en compte en conception ou en simulation, et affectant le comportement mécanique, électronique ou logiciel.
\end{itemize}

\subsubsubsection*{Répartition des \acrshort{MDF} par famille de sous-systèmes}

Les familles de sous-systèmes présentées précédemment sont principalement concernées par les modes de défaillance suivants :\\
\begin{itemize}
    \item \textbf{Éléments mécaniques statiques} : \acrshort{MDF}1, \acrshort{MDF}2
    \item \textbf{Éléments mécaniques dynamiques} : \acrshort{MDF}3, \acrshort{MDF}4
    \item \textbf{Éléments électroniques} : \acrshort{MDF}5, \acrshort{MDF}6, \acrshort{MDF}7
    \item \textbf{Éléments logiciels} : \acrshort{MDF}8, \acrshort{MDF}9, \acrshort{MDF}11
    \item \textbf{Facteurs humains (procéduraux)} : \acrshort{MDF}10, \acrshort{MDF}11
    \item \textbf{Facteurs externes (environnementaux)} : \acrshort{MDF}12 (et contribution possible à \acrshort{MDF}5,
    \acrshort{MDF}6, \acrshort{MDF}7)\\
\end{itemize}

Chacun de ces \acrlong{MDF} peut, selon le sous-système affecté et la phase de vol considérée, conduire à un ou plusieurs
\acrshort{ERP}.

\newpage
